\section{Method}\label{sec:method}

Existing schedulers treat spills as a local consequence after capacity is
exceeded and then apply a victim rule to a fixed order. Our experiments show
that this is the wrong primary lever: on a fixed order, four
Belady-style victim variants usually differ by at most $4\%$ in extra traffic,
whereas legal schedule orders can vary by more than an order of magnitude
(Section~\ref{sec:exp-mechanism}). Our method therefore optimizes schedule order
to shape the clean/dirty composition of overflow windows.

The method has three stages. First, spill diagnosis
(Section~\ref{sec:theory}) asks whether spills are unavoidable for the given DAG
and capacity. If they are, the objective shifts from eliminating spills to
lowering their cost, and overflow area becomes a cheap cross-stage surrogate.
Second, candidate-order generation and address assignment
(Section~\ref{sec:liveness-shaping}) produce complementary topological orders
and run best-fit placement with cost-aware spill insertion
(Algorithm~\ref{alg:assign-spill}). Third, we simulate the Cartesian product of
candidate orders and prefetch windows, then choose the best combination by a
lexicographic objective (Figure~\ref{fig:pipeline}).

\begin{figure*}[t]
  \centering
  \includegraphics[width=\textwidth]{\ksFigure{pipeline.pdf}}
  \caption{Method overview. The optimization framework consists of three stages:
  (1) inevitable spill determination to decide if optimization is needed;
  (2) three complementary candidate topological orders generated via liveness shaping;
  (3) Algorithm~\ref{alg:assign-spill} which executes address assignment and spill insertion along each order, followed by lexicographic key selection across candidate orders and prefetch windows.}
  \label{fig:pipeline}
\end{figure*}

\subsection{Spill Diagnosis and Surrogate Objective}\label{sec:theory}

If some schedule keeps every cache peak within capacity, then no spill is needed
and the clean/dirty distinction is irrelevant. The first question is therefore:
\emph{for this DAG and capacity, are spills unavoidable for all reasonable
schedules?} We use a linear-time certificate. Let $\mathcal{O}_\epsilon$ denote
schedules whose extra traffic is within $(1+\epsilon)$ of optimum, with
$\epsilon=0.1$. Let
$\underline{W}_c^\epsilon=\min_{S\in\mathcal{O}_\epsilon}\max_t W_c^S(t)$ be
the best per-cache peak working set within that near-optimal set.

\begin{theorem}[Spill-inevitability certificate]\label{thm:spill-inev}
If $\underline{W}_c^\epsilon>\Capc_c$, then every schedule in
$\mathcal{O}_\epsilon$ must spill on cache $c$. For any feasible implementation
$P$ of a schedule $S$, the extra traffic satisfies
$\extra(P)\ge\max_t(W_c^S(t)-\Capc_c)_+$. The proof is in the supplementary
material.
\end{theorem}

This certificate is deliberately one-sided: $\underline{W}_c^\epsilon>\Capc_c$
confirms that spills are unavoidable, but $\underline{W}_c^\epsilon\le\Capc_c$
does not prove that a spill-free implementation exists. We use it as an
applicability diagnostic rather than a complete feasibility test.

Peak live bytes are a direct way to compare candidate orders, but a P1-optimal
order can exceed capacity severely on a particular cache and produce high P2
traffic. We therefore use overflow area as the scheduling surrogate:
\begin{equation}\label{eq:area}
  A(S)=\sum_t\sum_{c\in\mathcal{C}}\bigl(\livec_c(t)-\Capc_c\bigr)_+ .
\end{equation}
In the implementation we use the normalized score
$\Phi(S)=\sum_t\sum_c(\livec_c(t)-\Capc_c)_+/\Capc_c$. The next result connects
this cheap surrogate to optimal spill traffic under bounded absence durations.

\begin{theorem}[Conditional overflow-area approximation]\label{thm:area-extra}
Assume every spilled interval has absence duration at most $L$. Suppose there
exists a feasible plan whose removed area satisfies $R(P)\le\lambda A(S)$ and
whose spilled intervals have absence duration at least $\ell$. If unit costs
obey $\kappa\in[1,2]$, then
\begin{equation}\label{eq:area-extra-bound}
  \frac{1}{L}\,A(S)\;\le\; E^\star(S)\;\le\; \frac{2\lambda}{\ell}\,A(S).
\end{equation}
The proof and a matching counterexample are in the supplementary material.
\end{theorem}

The surrogate $\Phi$ is useful for three reasons: it costs $O(N)$ to compute for
each candidate order, it gives a single scalar used consistently across P1, P2,
and P3, and Theorem~\ref{thm:area-extra} supplies a theoretical explanation for
why it should track spill traffic in capacity-bound regimes. It is not a global
optimality guarantee for schedule selection.

\subsection{Candidate Orders and Spill Insertion}\label{sec:liveness-shaping}

The candidate-order design is structure-agnostic. Orders depend only on buffer
clean/dirty state, size, and cache capacity; they do not encode operator-specific
motifs. The candidates share a memory-aware list-scheduling skeleton with three
ready sets: frees, operations, and allocations. Ready frees are scheduled first
to reduce residency pressure. Ready allocations are delayed until all ready
frees and operations have been scheduled, which compresses allocation-to-free
distance. We instantiate three complementary orders:

\begin{enumerate}[leftmargin=2em,label=(\arabic*)]
  \item \textbf{Pressure-aware order.} Allocation nodes are ranked by successor
  indegree. A smaller indegree indicates that consumers will become ready soon,
  so the allocation can be delayed until it is close to use.

  \item \textbf{Capacity-throttled order.} Allocation nodes are ranked by
  identifier and overlaid with capacity gating: an allocation that would push a
  cache over capacity is held back until no further delay is possible, with
  capacity weights normalized by $1/\Capc_c$ to balance caches of differing
  tightness.

  \item \textbf{ID-reserve order.} Within the three node classes, this order
  chooses the minimum identifier and performs no capacity throttling. It is a
  method candidate, not an external baseline: its purpose is to preserve the
  stable input-buffer order that often keeps clean \textsc{copy-in} buffers resident
  through compute windows, leaving cheap eviction reserve in overflow regions.
\end{enumerate}

Given a schedule order, Algorithm~\ref{alg:assign-spill} performs address
assignment and spill insertion along the order. When capacity is insufficient,
it selects an eviction victim by
\begin{equation}\label{eq:victim}
  \mathrm{victim}=\arg\max_b\; d(b)\cdot\frac{s_b}{e_b}
  =\arg\max_b\begin{cases}d(b), & b\text{ clean},\\ d(b)/2, & b\text{ dirty},\end{cases}
\end{equation}
where $d(b)$ is the next-use distance. Dirty buffers have their effective
distance halved, so the engine prefers cheap clean evictions when future-use
distances are comparable.

\begin{proposition}[Belady-margin stability]\label{prop:belady-margin}
Consider distance-dominant scoring rules whose cost-factor ratio satisfies
$\rho=g_{\max}/g_{\min}\le 2$. If the minimum next-use distance among selected
victims is greater than $\rho$ times the maximum next-use distance among
non-victims, then all such rules select the same victim set. The proof is in the
supplementary material.
\end{proposition}

Thus schedule order, not the victim rule, is the dominant axis on spill cost; we
quantify both effects in Section~\ref{sec:exp-mechanism}.

\begin{algorithm}[t]
\caption{Address assignment and spill insertion}\label{alg:assign-spill}
\begin{algorithmic}[1]
\Require Topological order $S$, capacities $\{\Capc_c\}$, prefetch window $H$
\Ensure Extended order, address assignment, spill list
\State Initialize a best-fit free-list manager $F_c$ for each cache $c$
\For{$(t,v)\in\mathrm{enumerate}(S)$}
  \If{$v=\textsc{Alloc}(b)$}
    \State $\mathrm{offset}\gets\Call{Place}{b,t}$
    \State $\mathrm{memory}[b]\gets\mathrm{offset}$
  \ElsIf{$v=\textsc{Free}(b)$}
    \If{$b$ is spilled} \Call{Reload}{$b,t$} \EndIf
    \State $F_{\tau(b)}.\mathrm{Release}(\mathrm{offset}_b,s_b)$
  \Else
    \For{$b\in v.\mathrm{bufs}$ where $b$ is spilled}
      \State \Call{Reload}{$b,t$}
    \EndFor
    \For{$b\in\mathrm{SpilledSet}$}
      \If{$d(b)-t\le H$ and $F_{\tau(b)}$ has $\ge s_b$ free}
        \State \Call{Reload}{$b,t$}
      \EndIf
    \EndFor
  \EndIf
\EndFor
\Statex
\Procedure{Place}{$b,t$}
  \State $\mathrm{offset}\gets F_{\tau(b)}.\mathrm{BestFit}(s_b)$
  \While{$\mathrm{offset}=\mathrm{nil}$}
    \State $v^*\gets\arg\max_{b'\in\mathrm{Resident}_{\tau(b)}} d(b')s_{b'}/e_{b'}$
    \State emit $\textsc{SpillOut}(v^*)$; release $(\mathrm{offset}_{v^*},s_{v^*})$
    \State $\mathrm{offset}\gets F_{\tau(b)}.\mathrm{BestFit}(s_b)$
  \EndWhile
  \State $F_{\tau(b)}.\mathrm{Allocate}(\mathrm{offset},s_b)$
  \State \Return $\mathrm{offset}$
\EndProcedure
\end{algorithmic}
\end{algorithm}

Reload timing does not change spill traffic, but it does affect whether reload
latency can overlap with computation. The prefetch window $H$ in
Algorithm~\ref{alg:assign-spill} controls this dimension: a spilled buffer whose
next use is within $H$ steps is reloaded early if free space is available and no
extra eviction is needed. For P2 we use $H\in\{0,5,40\}$ and select by
$(E,n,T)$. For P3 we use $H\in\{0,5,40,80,120\}$ and select by $(T,E,n)$. The
candidate set is monotone: adding a new order or prefetch window can only
improve or preserve the selected objective, because the final choice uses the
true P2/P3 key rather than the surrogate (see the supplementary material).

\paragraph{Complexity.}
The overall method runs in $O(K\cdot W\cdot(N\log N+E+B^2))$ time, where $K=3$ is
the number of candidate orders, $W$ is the prefetch-window grid size ($3$ for P2,
$5$ for P3), $N$ is the number of DAG nodes, and $B$ is the number of buffers.
Candidate-order generation costs $O(N\log N+E)$ (topological sort with a priority
queue). The bottleneck of address assignment is Belady victim selection: in the
worst case each allocation scans all resident buffers, giving $O(B)$ per
allocation over $B$ allocations, hence $O(B^2)$. The pipeline simulation is
$O(N+E)$ in the number of edges $E$; on the sparse DAGs of our benchmark this
term is below the combinatorial search and spill-assignment terms.
